\scanpage{27}
\uline{NB}\quad $\chi_\infty(E)=\varepsilon(\mathbb R f_*(E))-\varepsilon(\mathbb R f_!(E))$

Or on a vu que les deux $\varepsilon$ sont égaux, donc $\chi_\infty(E)=0$. De $=$ on a
\[
\delta_\infty(E)=\delta(\mathbb R f_*(E))\,\delta(\mathbb R f_!(E))^{-1}
\]
et on a vu que
\[
p^{(n+\rho)\chi(E)}\delta(\mathbb R f_*(E))^{-1}=\delta(\mathbb R f_!(E))
\]
d'où
\[
\left\{\begin{aligned}
\delta_\infty(E)&=p^{(n+\rho)\chi(E)}\delta(E)^{-2}
 =\deleted{\unsure{\delta}}\varepsilon(E)^{-2}\ \deleted{\unsure{=\delta'(E)q^{(n+\rho)\chi}|=\delta(E)p^{(n+\rho)\chi(E)}}}\\
\deleted{\unsure{\delta}}\varepsilon(E)&=\deleted{\unsure{\varepsilon}}\delta(E)\,p^{-(n+\rho)\chi(E)/2}
 \qquad (=\pm1\ \text{si $X$ est \deleted{\unsure{cpt}} propre})
\end{aligned}\right.
\]
(qui est bien $1$ lorsque $X$ est \unsure{compacte|compact|complet} !)

d'où l'équation fonctionnelle
\[
\boxed{L_E^\infty\left(\frac{1}{p^{n+\rho}t}\right)L_E^\infty(t)\,\delta_\infty(E)=1}
\]
\[
\left(\delta_\infty(E)=p^{(n+\rho)\chi(E)}\delta(E)^{-2}\deleted{\unsure{)}}=\varepsilon(E)^{-2}\right)
\]

$\bigl[$ Invariants arithmétiques équivalents, (moyennant la connaissance de l'invariant \uline{géométrique} $\chi(E)$, et de $\unsure{q|\varphi}=\deleted{\unsure{p_s}}p^{n+\rho}$)

\margin{\uline{NB.} $\varepsilon(E)=\pm1$ si $X$ \unsure{complet|complète}}
% slanted
\[
\delta_{.}(E)=\delta_!^*(E),\quad \delta'(E)=\delta_{.}^*(E),\quad \varepsilon(E),\quad \delta_\infty(E)\ \biggr]
\]
\[
\boxed{\delta(E)\delta'(E)=q^{\chi(E)}}
\]
\[
\boxed{\varepsilon(E)=\delta(E)q^{-\chi(E)/2}=\delta'(E)q^{-\chi(E)/2}}
\]
\margin{$\searrow$ \unsure{variance|variante} avec $E,\unsure{\lambda|n|X}$ !!}
\[
\boxed{\begin{aligned}
\delta_\infty(E)&=q^{\chi(E)}\delta(E)^{-2}=q^{-\chi(E)}\delta'(E)^2\\
&=\varepsilon(E)^{-2}
\end{aligned}}
\]
% archive sheet number: 26

